\section{Integer boxes and the winning-set reduction}
\label{sec:boxes}

The game will be solved by recovering one integer potential. We first
construct its operator and prove a comparison principle in an arbitrary
finite box. The principle gives both a unique clipped fixed point and an
exact, possibly slow, iteration. We then identify the winning set from the
fixed point in a particular wide box. This order separates the reusable
comparison result from the width needed for the game reduction.

Potential transformations originate in cyclic-game algorithms
\cite{GurvichKarzanovKhachiyan1988}; bounded monotone integer updates also
underlie energy progress measures
\cite[Definitions~1--2 and Lemmas~4--6]{BrimEtAl2011}.
The symmetric box statements below are proved directly.

\subsection{The transformed operator}

Use coordinatewise order, maxima, and minima throughout. A scalar added to
a vector is added in every coordinate. For integer vectors $a\le b$, set
\[
 [a,b]=\{p\in\Z^V:a_i\le p_i\le b_i\text{ for all }i\in V\},
 \qquad
 (\clip_{[a,b]}p)_i=\min\{b_i,\max\{a_i,p_i\}\}.
\]
A positive integer $B$ is a \emph{reference width} if $b_i-a_i\le B$ for
all $i$. Endpoints may be negative or vary with $i$, and an interval may
collapse to a single integer.

Apply the integer cycle perturbation used by Bj\"orklund and
Vorobyov~\cite[Section~9]{BjorklundVorobyov2007}, and define
\begin{equation}
 t(e)=(n+1)w(e)+1,\qquad
 W=\max\bigl\{1,\max_{e\in E}|t(e)|\bigr\},\qquad B_*=8nW.
 \label{eq:transformed}
\end{equation}
A simple directed cycle has no repeated vertex except its two endpoints;
a self-loop therefore has length one.

\begin{lemma}[Cycle signs]\label{box:cycle-signs}
Every simple directed cycle has nonzero $t$-sum. Its $t$-sum is positive
exactly when its $w$-sum is nonnegative, and negative exactly when its
$w$-sum is negative.
\end{lemma}

\begin{proof}
For a cycle of length $\ell\in\{1,\ldots,n\}$ and integer $w$-sum $c$,
the new sum is $(n+1)c+\ell$. If $c\ge0$ this is at least one. If $c<0$,
then $c\le-1$, and the new sum is at most $-(n+1)+n=-1$.
The calculation concerns the actual edges of the cycle and so also covers
parallel edges.
\end{proof}

On the whole space $\Z^V$, define
\begin{equation}\label{eq:operator}
 F_i(p)=
 \begin{cases}
  \displaystyle\max_{e:i\to j}(t(e)+p_j),&i\in V_{\mathrm{Max}},\\[2mm]
  \displaystyle\min_{e:i\to j}(t(e)+p_j),&i\in V_{\mathrm{Min}}.
 \end{cases}
\end{equation}
The graph is finite and has no sinks, so these extrema exist and are
attained. Two properties will be used repeatedly:
\begin{equation}\label{box:operator-properties}
 p\le q\ \Longrightarrow\ F(p)\le F(q),\qquad
 F(p+c)=F(p)+c\quad(c\in\Z).
\end{equation}
Defining $F$ outside any particular box is essential: later comparisons
will be translated before being clipped into new boxes.

\subsection{Comparison and exact iteration}

\begin{definition}[Box comparisons]\label{def:box-comparisons}
For a finite integer box $[a,b]$, a \emph{subsolution} is a full vector
$y\in[a,b]$ satisfying
\[
 y_i\le F_i(y)\quad(i\in S_y),\qquad S_y=\{i\in V:y_i>a_i\}.
\]
A \emph{supersolution} is a full vector $z\in[a,b]$ satisfying
\[
 z_i\ge F_i(z)\quad(i\in S_z),\qquad S_z=\{i\in V:z_i<b_i\}.
\]
We call $S_y$ and $S_z$ their \emph{supports}, always relative to the
specified box. At a collapsed coordinate $a_i=b_i$, neither support
contains $i$ and neither inequality is required.
\end{definition}

The comparison inequalities are imposed only on their respective
supports. Nevertheless, their conclusion orders the vectors everywhere.

\begin{lemma}[Comparison principle]\label{lem:box-comparison}
Let $F$ be the transformed operator in Equation~\eqref{eq:operator}.
For arbitrary integer bounds $a\le b$, every subsolution $y\in[a,b]$
and supersolution $z\in[a,b]$ satisfy $y_i\le z_i$ for \emph{all}
$i\in V$. This includes boxes with negative, nonconstant, or collapsed
coordinate intervals.
\end{lemma}

\begin{proof}
Suppose that $d=\max_i(y_i-z_i)>0$, and take a vertex attaining this
maximum. There $y_i>z_i\ge a_i$ and $z_i<y_i\le b_i$, so both defining
inequalities apply. Since $y\le z+d$, monotonicity and translation give
\[
 d=y_i-z_i\le F_i(y)-F_i(z)\le d.
\]
Equality forces both comparison slacks to vanish:
\begin{equation}\label{box:comparison-equalities}
 y_i=F_i(y),\qquad z_i=F_i(z).
\end{equation}
We claim that a common tight edge leads to another vertex attaining $d$.
At a Max vertex choose an edge $e:i\to j$ attaining $F_i(y)$; at a Min
vertex choose one attaining $F_i(z)$. In either case,
\[
 d=F_i(y)-F_i(z)
 \le (t(e)+y_j)-(t(e)+z_j)=y_j-z_j\le d.
\]
Every inequality is therefore an equality. The edge is tight for both
vectors, and its head again attains the maximum difference.

Continue choosing edges in this way. Finiteness supplies a repeated
vertex; stopping at the first repetition produces a simple cycle of
edges satisfying $y_i=t(e)+y_j$. Summing these equalities gives a zero
$t$-sum, in conflict with Lemma~\ref{box:cycle-signs}. Thus no positive
maximum difference exists, proving the assertion on all coordinates.
\end{proof}

\begin{lemma}[The box solution and basic iteration]\label{lem:box-solution}
For the same operator $F$ and any finite integer box $[a,b]$, there is
exactly one vector $x\in[a,b]$ that is both a subsolution and a
supersolution. It is exactly the unique fixed point of
\[
 T_{a,b}(p)=\clip_{[a,b]}F(p).
\]
The synchronous procedure $\Basic(a,b)$ initializes $p=a$, evaluates
$q=T_{a,b}(p)$, returns $p$ if $q=p$, and otherwise replaces $p$ by $q$
and repeats. It returns $x$ after at most
\[
 1+\sum_{i\in V}(b_i-a_i)\le1+nB
\]
operator evaluations, including the final stability test, for any
reference width $B$. No positivity or nondegeneracy of the coordinate
intervals is assumed.
\end{lemma}

\begin{proof}
First inspect the fixed-point condition in one coordinate. In the strict
interior it says $F_i(p)=p_i$, the conjunction of the two comparison
inequalities. At $p_i=a_i<b_i$, clipping fixes $p_i$ exactly when
$F_i(p)\le p_i$, the active supersolution inequality. At
$a_i<p_i=b_i$, it fixes $p_i$ exactly when $F_i(p)\ge p_i$, the active
subsolution inequality. A collapsed coordinate is always fixed by
clipping and imposes neither comparison. Fixed points are thus precisely
the vectors having both comparison properties.

The map $T_{a,b}$ is monotone and takes the box into itself. Its first
iterate is at least $a$, so all successive iterates from $a$ are
nondecreasing. Each strict update increases the integer sum of the
coordinates by at least one, whereas the available total increase is
$\sum_i(b_i-a_i)$. This proves existence, termination, and the count,
with one extra evaluation to recognize stability. If $x,x'$ are fixed
points, apply Lemma~\ref{lem:box-comparison} to them in both orders to
obtain $x\le x'$ and $x'\le x$. Hence $x=x'$.
\end{proof}

We call $x$ the \emph{box solution}. The count above may be enormous for
binary-encoded endpoints. Our recursive solver will use $\Basic$ only
at small reference width; the next result explains why recovering the
solution of one much wider box is sufficient.

\subsection{The two potential bands}

\begin{proposition}[Reduction to a box solution]\label{prop:game-reduction}
Let $x$ solve the constant box $[0,B_*]$, with $B_*=8nW$ as in
Equation~\eqref{eq:transformed}. Then the original game's exact
zero-threshold winning set is
\[
 U=\{i\in V:x_i>B_*/2\}.
\]
A positional Max strategy guarantees original mean payoff at least zero
from every vertex of $U$. Outside $U$, a positional Min strategy ensures
that every consistent play has original upper limiting average at most
$-1/n$, including against arbitrary history-dependent Max strategies.
\end{proposition}

\begin{proof}
We first locate all coordinates of $x$. At an interior coordinate,
$x_i=F_i(x)$, so an attaining edge is tight: $x_i=t(e)+x_j$. Follow
such edges while the vertices remain interior. A repetition would give
a simple zero-sum $t$-cycle, which Lemma~\ref{box:cycle-signs} excludes.
A boundary is therefore reached within $n-1$ edges. For a coordinate
already on a boundary the path has length zero. Telescoping proves that
every coordinate lies within $nW$ of one boundary, giving the partition
\begin{equation}\label{box:bands}
 V_{\mathrm{lo}}=\{i:0\le x_i\le nW\},\qquad
 V_{\mathrm{hi}}=\{i:7nW\le x_i\le8nW\}.
\end{equation}
The midpoint $B_*/2=4nW$ separates these bands.

On the high band the subsolution inequality is active. At a high Max
vertex choose an attaining edge; at a high Min vertex keep every edge.
All these permitted edges satisfy
\begin{equation}\label{box:high-edge}
 x_i\le t(e)+x_j.
\end{equation}
Their heads must also be high, since a low head would imply
\[
 7nW\le x_i\le t(e)+x_j\le(n+1)W<7nW
\]
for $n\ge1$. The selected Max edges, extended arbitrarily at low
vertices, define a positional strategy preserving the high band against
all choices of Min.

Every simple cycle in this permitted graph has nonnegative $t$-sum by
Equation~\eqref{box:high-edge}; hence it has nonnegative $w$-sum by
Lemma~\ref{box:cycle-signs}. Set $M=\max_e|w(e)|$. Decompose any
length-$T$ prefix of any consistent play into simple cycles and a
residual walk: repeatedly remove the segment ending at the first
repeated vertex. The residual walk has at most $n-1$ edges, while the
removed cycles have nonnegative original sums. The prefix weight is at
least $-(n-1)M$. Dividing by $T$ and taking the lower limit proves
Max's assertion, with no restriction on Min's dependence on the past.

On the low band the supersolution inequality is active. At a low Min
vertex choose an attaining edge, and at a low Max vertex keep all edges.
These edges satisfy
\begin{equation}\label{box:low-edge}
 x_i\ge t(e)+x_j.
\end{equation}
A high head would give
\[
 nW\ge x_i\ge t(e)+x_j\ge(7n-1)W>nW,
\]
so this permitted graph stays low. Extend the chosen Min edges to a
positional strategy on all vertices. Each permitted simple cycle has
nonpositive $t$-sum and therefore strictly negative integer $w$-sum,
at most $-1$. In a decomposition of a length-$T$ prefix into $r$ simple
cycles and a residual walk, the cycle lengths are at most $n$ and the
residual length at most $n-1$. Consequently
\[
 r\ge\frac{T-(n-1)}n,
 \qquad
 \sum_{j<T}w(e_j)\le-r+(n-1)M
 \le-\frac{T-(n-1)}n+(n-1)M.
\]
The upper limiting average is at most $-1/n$. This estimate holds for
every play allowed by the chosen Min strategy, including those from
history-dependent Max play. The partition into the two bands now proves
the exact winning-set formula.
\end{proof}
